%=====================================================================
\chapter{Responsible AI and Its Limits}\label{ch:responsible}
%=====================================================================
\locator{6}{Part VI --- Learning, Language and Responsible AI}

\begin{outcomes}
By the end of this chapter you will be able to:
\begin{tight}
  \item \textbf{Identify} the points at which bias enters a rule-based system,
        none of which require anyone to intend it.
  \item \textbf{Audit} a deployed rule set for differential treatment, and
        \textbf{present} the evidence.
  \item \textbf{Argue} why an explanation is an accountability artefact rather than
        a convenience.
  \item \textbf{State} the computational limits of the methods in this book, with
        numbers.
  \item \textbf{Recognise} brittleness at the edge of competence, and
        \textbf{distinguish} it from ordinary error.
\end{tight}
\end{outcomes}

\begin{prereq}
\chref{expert} (the rule engine audited here), \chref{uninformed} (the branching
factor that sets the computational limit) and \chref{trends} (the contrast with
learned systems). This chapter takes the systems you built and asks what they owe
the people they are used on.
\end{prereq}

\begin{keyterms}
bias $\bullet$ proxy variable $\bullet$ differential treatment $\bullet$ disparate
impact $\bullet$ audit $\bullet$ explainability $\bullet$ accountability $\bullet$
contestability $\bullet$ automation bias $\bullet$ brittleness $\bullet$ graceful
degradation $\bullet$ intractability $\bullet$ the limits of competence
\end{keyterms}

\section{Nobody Intended It}

\chref{expert}'s escalation adviser is twenty rules about severity, customer tier,
workarounds and on-call budgets. There is no rule about teams, no rule about people,
and nobody who wrote it intended to treat anybody differently.

Run it over two hundred tickets and it escalates $43\%$ of the ones reported by one
team and $21\%$ of those reported by another.

That is this chapter. Not malice, not a bug in the ordinary sense, and not something
that inspection of the rules would have revealed --- a property of the system that
only appears when you measure the outcomes.

\section{Where Bias Enters a Rule Base}

\dsfig{%
\begin{tikzpicture}[font=\scriptsize,
  s/.style={rounded corners=3pt, draw=#1, line width=1.1pt, fill=#1!8,
            text=#1!50!black, text width=3.15cm, align=flush left,
            inner sep=6pt, minimum height=3.1cm},
  hd/.style={rounded corners=2pt, fill=#1, text=white,
             font=\tiny\bfseries, inner sep=2.5pt, text width=3.0cm,
             align=center},
  a/.style={-{Stealth[length=2mm]}, neutral!70, line width=1pt}]

\node[s=primary]   (b1) at (0,0)     {\\[2pt]A rule uses a feature that
  \emph{stands in} for something you did not mean to use.\\[4pt]
  \emph{Here:} \texttt{workaround = yes} is recorded far more often by teams with
  mature runbooks.};
\node[s=goldhl]    (b2) at (3.85,0)  {\\[2pt]The facts fed to the system are
  produced unevenly.\\[4pt]
  \emph{Here:} an experienced reporter knows to set severity to \emph{major}; a new
  one leaves the default.};
\node[s=accent]    (b3) at (7.70,0)  {\\[2pt]The rules are correct and their
  \emph{combination} is not.\\[4pt]
  \emph{Here:} R3 and R7 are each defensible; together they route an entire team's
  tickets to the queue.};
\node[s=tealhl]    (b4) at (11.55,0) {\\[2pt]The measure itself encodes a
  choice.\\[4pt]
  \emph{Here:} counting \emph{breaches avoided} and not \emph{who waited} makes one
  group's waiting invisible.};

\node[hd=primary] at (b1.north) {1. PROXY VARIABLES};
\node[hd=goldhl]  at (b2.north) {2. UNEVEN INPUTS};
\node[hd=accent]  at (b3.north) {3. INTERACTION};
\node[hd=tealhl]  at (b4.north) {4. THE MEASURE};

\node[rounded corners=3pt, fill=lightbg, draw=primary!30, line width=0.9pt,
      text=primary!70!black, font=\scriptsize, text width=14.2cm,
      align=flush left, inner sep=6pt] at (5.8,-2.3)
  {\textbf{None of the four requires anyone to intend anything}, and none is
   visible in the rules. You cannot find any of them by reading the rule base, and
   three of the four survive deleting every sensitive attribute from the inputs.
   They are found by measuring outcomes across groups --- which means the audit is
   not an optional extra but the only instrument that detects the problem.};
\end{tikzpicture}}%
{Four ways bias enters a system whose rules mention nothing sensitive. The first
--- a \term{proxy variable} --- is the one that survives every attempt to fix the
problem by deleting columns.}{bias-entry}

\begin{alertbox}[Deleting the sensitive attribute does not work]
The instinctive fix is to remove the team name from the inputs. The escalation
adviser never had it.

Proxies are the reason. \texttt{workaround = yes} correlates with team maturity;
severity correlates with reporter experience; both are legitimate operational facts
that a rule ought to use. Remove them and the system gets worse at its job without
becoming fairer, because the correlation lives in the world rather than in the
column.

This is why \emph{fairness through unawareness} --- not looking --- is the one
approach that reliably fails. The alternative is not to look less but to measure
more: keep the sensitive attribute \emph{for the audit}, and never for the decision.
\end{alertbox}

\section{Auditing a Deployed Rule Set}

An audit is three steps and about thirty lines of code.

\begin{enumerate}[leftmargin=2.2em, itemsep=3pt, topsep=4pt]
  \item \textbf{Measure the outcome rate per group.} Run the deployed system over a
        representative population and count.
  \item \textbf{Find the rule responsible.} For each rule, ask what the gap would be
        if that rule were removed. The one whose removal closes the gap most is the
        one to examine.
  \item \textbf{Read that rule's explanation chain for an affected case}, and decide
        whether the reasoning is defensible. This is a human judgement and the
        chain is what makes it possible.
\end{enumerate}

\begin{examplebox}[The audit, on two hundred tickets]
\begin{tabular}{@{}L{4.2cm}L{2.2cm}L{2.4cm}L{4.8cm}@{}}
\toprule
\textbf{Reporting team} & \textbf{Tickets} & \textbf{Escalated} &
\textbf{Rate} \\
\midrule
Platform (mature runbooks) & 100 & 21 & $21\%$ \\
Payments (newer team) & 100 & 43 & $43\%$ \\
\midrule
\textbf{Gap} & & & \textbf{22 percentage points} \\
\bottomrule
\end{tabular}

\smallskip
Step 2 points at the pathway \textbf{R3} $\to$ \textbf{R7}: a major ticket from a
standard-tier customer becomes P2 (R3), and a P2 with a workaround needs no owner
(R7), so it is never escalated. Neither rule is wrong; both are entirely sensible.
The disparity is that \texttt{workaround} is recorded far more often by the team
with mature runbooks, so that pathway fires overwhelmingly for one group.
\end{examplebox}

\begin{alertbox}[Ablation has a trap, and the lab walks into it deliberately]
Asked simply ``which removal closes the gap most?'', the ablation answers
\textbf{R9} --- and closes the gap completely, to zero points.

R9 is the only rule that concludes \texttt{escalate = yes}. Removing it means
nothing is ever escalated, by anyone, so the rates are $0\%$ and $0\%$ and the
disparity is perfect equality. The same happens with R5, which is upstream of it.

The measure was satisfied by destroying the outcome. That is a general hazard of
optimising a fairness metric rather than understanding a system, and the guard is
simple: \emph{report the overall rate alongside the gap}, and discard any ablation
that collapses it. Among the rules that leave the outcome intact, R3 carries the
gap --- $22$ points down to $9$.
\end{alertbox}

\begin{conceptbox}
What should be done about it is not a technical question, and the audit's value is
that it makes the choice explicit rather than making it by default.

Three defensible responses. \emph{Accept it}: the teams genuinely differ in
maturity, the escalations reflect real need, and the gap is the system working. Fix
the input: help the newer team write runbooks, so the difference in what is recorded
reflects a difference in the world rather than in documentation practice. Or
\emph{change the rule}: require an owner for P2 tickets from teams without an
established runbook, regardless of the workaround field.

All three are reasonable. What is not reasonable is not knowing --- which was the
situation before the audit, and which is the default state of every deployed rule
base that nobody has measured.
\end{conceptbox}

\section{Explanation as Accountability}

\chref{expert} produced an explanation chain and called it the deliverable. Here is
the argument for that claim.

An explanation is not primarily for the user's comfort. It is the artefact that
makes a decision \term{contestable}: it tells a person exactly which facts and which
rules produced an outcome, so they can dispute a \emph{specific} step rather than
the conclusion in general. ``The system says no'' cannot be argued with; ``R7 fired
because \texttt{workaround} was recorded as yes'' can, and the argument is about
something checkable.

\smallskip
\noindent\begin{longtable}{@{}L{3.4cm}L{5.5cm}L{5.6cm}@{}}
\toprule
\rowcolor{primary!15}
\textbf{Question} & \textbf{Rule engine (\chref{expert})} &
\textbf{Learned model (\chref{trends})} \\
\midrule
\endhead
Why this outcome? & The chain of rules and facts, checkable line by line &
A post-hoc rationale, generated the same way as the decision \\
What would change it? & Read the chain: change any fact in it & Unclear without
probing the model \\
Who is responsible? & Whoever wrote the rule that fired --- a named person &
Diffuse: the data, the objective, the training \\
Can a subject contest it? & Yes, on a specific step & Only in general \\
Can it be corrected? & Edit one rule & Retrain, and hope \\
\bottomrule
\caption{Why explanation is an accountability property. The last two rows are the
ones that matter to somebody affected by the decision.}\label{tab:explain}
\end{longtable}

\begin{alertbox}[Automation bias]
There is a failure that belongs to the people rather than the system.

Given a system that is usually right, human reviewers stop reviewing. They accept
the recommendation, including in the cases where they would have known better ---
and the effect is strongest exactly where the system is most often correct, because
that is where the habit of deferring forms.

An explanation helps only if somebody reads it, and an explanation that is always
present and always plausible becomes wallpaper. The practical mitigations are
procedural rather than technical: sample decisions for genuine review, measure the
override rate and treat a rate near zero as a warning rather than a success, and
make the system say when it is near the edge of its competence --- which is the
subject of the next section.
\end{alertbox}

\section{Brittleness and the Edge of Competence}

\chref{expert} named brittleness as one of the two causes of the expert-system
collapse. It is worth being precise about what it means, because it is not the same
as being wrong sometimes.

A person with expertise, asked something outside it, says so. A rule base asked
something outside its competence either fires nothing --- unhelpful but honest ---
or fires a rule whose conditions technically match and whose conclusion is absurd,
with exactly the same confidence as for a routine case. There is no representation
of its own boundary, so there is no signal.

\term{Graceful degradation} is the property it lacks: performance that falls off
smoothly as inputs move away from the familiar, and --- more importantly --- a
measure that falls with it.

\begin{conceptbox}
Learned systems degrade more smoothly and \emph{also} give no reliable signal: a
language model is as fluent about things it has no basis for as about things it
does. So neither tradition solves this, and the difference is in what you can do
about it.

With a rule base you can at least \emph{characterise} the boundary, because the
knowledge is explicit: enumerate the situations no rule covers, and say so. That is
a real and underused capability, it takes an afternoon, and the lab does it.
\end{conceptbox}

\section{Computational Limits}

The last limit is arithmetic, and it applies to every method in Parts~II to~V.

\smallskip
\noindent\begin{longtable}{@{}L{4.2cm}L{4.6cm}L{5.7cm}@{}}
\toprule
\rowcolor{primary!15}
\textbf{Method} & \textbf{Cost} & \textbf{What that means in practice} \\
\midrule
\endhead
Uninformed search (\chref{uninformed}) & $O(b^{d})$ & At $b=4$, depth $10$ is a
million nodes and depth $20$ is $10^{12}$ \\
$A^{*}$ (\chref{informed}) & Exponential in general & A good heuristic moves the
constant, not the exponent \\
Constraint satisfaction (\chref{csp}) & NP-complete & Solvable in practice by
propagation; no guarantee \\
Adversarial search (\chref{adversarial}) & $O(b^{d})$, $O(b^{d/2})$ with perfect
ordering & Twice the depth for the same time, and still exponential \\
Propositional entailment (\chref{propositional}) & co-NP-complete & $2^{n}$ models;
DPLL is fast on structured instances \\
First-order entailment (\chref{fol}) & Semi-decidable & No proof may mean no proof
exists, or may mean not yet \\
Bayesian inference (\chref{bayesnets}) & NP-hard in general & Linear on a chain,
exponential on a dense graph \\
\bottomrule
\caption{The computational limits of this course's methods. Every entry is
exponential, NP-hard or undecidable, and no engineering removes
that.}\label{tab:limits}
\end{longtable}

The practical consequence is not despair. Every one of these methods is used daily
on instances that matter, because real instances have structure --- and the whole
content of Parts~II to~V is techniques for exploiting it: heuristics, propagation,
pruning, conditional independence, unit propagation. What the table forbids is the
belief that a faster machine will do instead. At $b = 4$, a hundredfold speedup buys
between three and four extra levels of depth, and then stops.

\begin{worked}[Auditing the escalation adviser]
Two hundred tickets, evenly split between two reporting teams. The rule base is
\chref{expert}'s, unchanged, and it mentions no team.

\smallskip
\textbf{Step 1: measure.} Escalation rates are $21\%$ for Platform and $43\%$ for
Payments --- a gap of $22$ points, against an overall rate of $32\%$. Nothing in
the rules explains it.

\smallskip
\textbf{Step 2: attribute.} Remove one rule at a time and re-measure both the gap
\emph{and} the overall rate:
\begin{tight}
  \item Without R9: the gap is $0$ --- and so is the escalation rate. R9 is the only
        rule concluding \texttt{escalate = yes}, so its removal achieves equality by
        abolishing the outcome. Discard it, and R5 with it.
  \item Without R3 (\emph{major and standard tier gives P2}): the gap falls from
        $22$ to $9$ points, and the overall rate stays near $32\%$. This is the
        real answer.
  \item Without any other single rule: the gap barely moves.
\end{tight}
So the R3 $\to$ R7 pathway carries most of it. The technique is \emph{ablation}; it
required no insight, only the discipline to try each rule and to report the overall
rate beside the gap.

\smallskip
\textbf{Step 3: examine.} Take a Platform ticket that R7 kept out of the queue and
read its chain: \texttt{escalate = no} by R11, because \texttt{needs\_owner = no} by
R7, because the priority was P2 and \texttt{workaround = yes}.

Now ask the operational question: \emph{why is \texttt{workaround} recorded more
often for Platform?} The answer is that Platform has runbooks and Payments does not,
so Platform's reporters have a documented workaround to cite and Payments' reporters
leave the field empty.

\smallskip
\textbf{What the audit established.} The system is doing exactly what it was told.
R7 is a good rule. The disparity comes from a difference in \emph{documentation
practice} being read as a difference in \emph{operational need}.

\smallskip
\textbf{And what it did not establish.} Whether that is acceptable. If Platform
genuinely needs fewer escalations because it has runbooks, the gap is the system
working. If Payments is being under-served because nobody has written their runbooks
yet, it is not. \emph{That is a management decision}, and the audit's contribution
is to place it in front of a person who can make it, with the evidence attached ---
rather than leaving it to be made silently, by a rule nobody re-read, on two hundred
tickets a week.
\end{worked}

\begin{pitfall}
\begin{tight}
  \item \textbf{Fairness through unawareness.} Deleting the sensitive attribute
        leaves the proxies. The escalation adviser never had the team name and the
        gap was $22$ points.
  \item \textbf{Auditing the rules instead of the outcomes.} Every rule here is
        defensible in isolation. Only measurement across groups finds the problem.
  \item \textbf{Treating an explanation as a user-experience feature.} It is what
        makes a decision contestable and correctable, and those are properties an
        affected person needs, not a designer.
  \item \textbf{Reading a low override rate as success.} It is the signature of
        automation bias. Sample for genuine review and watch the rate.
  \item \textbf{Confusing brittleness with ordinary error.} Ordinary error is being
        wrong sometimes; brittleness is being wrong with undiminished confidence
        just outside the region of competence, with no signal.
  \item \textbf{Believing hardware solves intractability.} At $b=4$ a hundredfold
        speedup buys three or four levels of depth. Structure is the only real
        answer, which is what Parts~II to~V were about.
\end{tight}
\end{pitfall}

%---------------------------------------------------------------------
\begin{lab}[Auditing the Escalation Rules]
The rule base of \chref{expert}, run over two hundred synthetic tickets, measured by
group, attributed by ablation, and explained. Then the competence boundary and the
computational limit.
\begin{lstlisting}[language=Python]
"""Chapter 21 lab -- auditing a rule base nobody wrote badly.

The rules mention no team. The escalation rate differs by 38
percentage points between two teams. Everything here is measurement:
the disparity is invisible in the rule base and obvious in the
outcomes.
"""

import random

# ---- Chapter 13's rule base, unchanged ----------------------------
RULES = [
    ("R1",  [("severity", "critical")], ("priority", "P1")),
    ("R2",  [("severity", "major"), ("tier", "platinum")],
     ("priority", "P1")),
    ("R3",  [("severity", "major"), ("tier", "standard")],
     ("priority", "P2")),
    ("R4",  [("severity", "minor")], ("priority", "P3")),
    ("R5",  [("priority", "P1")], ("needs_owner", "yes")),
    ("R6",  [("priority", "P2"), ("workaround", "no")],
     ("needs_owner", "yes")),
    ("R7",  [("priority", "P2"), ("workaround", "yes")],
     ("needs_owner", "no")),
    ("R8",  [("priority", "P3")], ("needs_owner", "no")),
    ("R9",  [("needs_owner", "yes"), ("owner_available", "no")],
     ("escalate", "yes")),
    ("R11", [("needs_owner", "no")], ("escalate", "no")),
]


def forward_chain(facts, rules=RULES):
    wm, fired = dict(facts), set()
    while True:
        applicable = [r for r in rules
                      if r[0] not in fired
                      and all(wm.get(a) == v for (a, v) in r[1])
                      and r[2][0] not in wm]
        if not applicable:
            return wm, fired
        applicable.sort(key=lambda r: (-len(r[1]), rules.index(r)))
        r = applicable[0]
        wm[r[2][0]] = r[2][1]
        fired.add(r[0])


def explain(goal, facts, rules=RULES, seen=None):
    """Backward chaining, for the chain an auditor reads."""
    seen = seen or set()
    if facts.get(goal[0]) == goal[1]:
        return ("fact", goal)
    if goal in seen:
        return None
    for r in rules:
        if r[2] != goal:
            continue
        subs = []
        for c in r[1]:
            got = explain(c, facts, rules, seen | {goal})
            if got is None:
                subs = None
                break
            subs.append(got)
        if subs is not None:
            return ("rule", r[0], goal, subs)
    return None


def render(tree, indent=0):
    pad = "  " * indent
    if tree[0] == "fact":
        return ["%s%s = %s  (given)" % (pad, tree[1][0], tree[1][1])]
    _, name, goal, subs = tree
    out = ["%s%s = %s  (by %s, because)" % (pad, goal[0], goal[1], name)]
    for s in subs:
        out += render(s, indent + 1)
    return out


# ---- two hundred tickets from two teams ---------------------------
def population(n=200, seed=0):
    """The ONLY difference between the teams is how often a
    workaround is documented. Nothing else, and no rule sees the
    team name."""
    rng = random.Random(seed)
    out = []
    for i in range(n):
        team = "Platform" if i % 2 == 0 else "Payments"
        # Platform has mature runbooks, so a workaround is usually
        # recorded; Payments is a newer team and usually has none
        p_workaround = 0.80 if team == "Platform" else 0.15
        out.append({
            "id": "T-%04d" % (4000 + i),
            "team": team,                       # for the AUDIT only
            "severity": rng.choice(["critical", "major", "major",
                                    "minor"]),
            "tier": rng.choice(["platinum", "standard", "standard"]),
            "workaround": "yes" if rng.random() < p_workaround else "no",
            "owner_available": rng.choice(["yes", "no", "no"]),
        })
    return out


def decide(ticket, rules=RULES):
    facts = {k: v for k, v in ticket.items()
             if k not in ("id", "team")}        # the team is NOT input
    wm, fired = forward_chain(facts, rules)
    return wm.get("escalate", "no"), fired, facts


def rates(tickets, rules=RULES):
    counts = {}
    for t in tickets:
        esc, _, _ = decide(t, rules)
        a, b = counts.get(t["team"], (0, 0))
        counts[t["team"]] = (a + (esc == "yes"), b + 1)
    return {k: (a, b, 100.0 * a / b) for k, (a, b) in counts.items()}


if __name__ == "__main__":
    tickets = population()
    print("no rule mentions the reporting team. Rule attributes used:")
    print("   ", sorted({a for r in RULES for (a, _) in r[1]}))
    print()

    r = rates(tickets)
    print("=== step 1: measure the outcomes ===")
    print("   team        tickets  escalated   rate")
    print("   " + "-" * 40)
    for team in sorted(r):
        esc, tot, pct = r[team]
        print("   %-10s %8d %10d %6.0f%%" % (team, tot, esc, pct))
    gap = abs(r["Payments"][2] - r["Platform"][2])
    print("   gap: %.0f percentage points" % gap)
    assert gap > 15

    # --- step 2: which rule carries the gap? -----------------------
    overall = sum(x[0] for x in r.values()) / sum(x[1] for x in r.values())
    print("   overall escalation rate: %.0f%%" % (100 * overall))

    # --- step 2: which rule carries the gap? -----------------------
    print()
    print("=== step 2: ablation -- remove one rule at a time ===")
    print("   Report the OVERALL rate as well as the gap. A rule whose")
    print("   removal closes the gap by destroying the outcome has")
    print("   told you nothing.")
    print()
    print("   rule removed    gap   overall rate")
    print("   " + "-" * 38)
    scores = {}
    for victim in [x[0] for x in RULES]:
        reduced = [x for x in RULES if x[0] != victim]
        rr = rates(tickets, reduced)
        g = abs(rr["Payments"][2] - rr["Platform"][2])
        ov = (sum(x[0] for x in rr.values())
              / sum(x[1] for x in rr.values()))
        scores[victim] = (g, ov)
        flag = ""
        if ov < 0.5 * overall:
            flag = "  <- outcome collapsed; tells us nothing"
        print("   %-14s %5.0f %11.0f%%%s" % (victim, g, 100 * ov, flag))

    # R9 is the only rule concluding escalate=yes, so removing it
    # sets every rate to zero and the gap to zero. That is the trap.
    assert scores["R9"][0] == 0 and scores["R9"][1] == 0

    # the honest question: among rules that leave the outcome intact,
    # which one carries the gap?
    intact = {k: v for k, v in scores.items()
              if v[1] >= 0.5 * overall}
    worst = min(intact, key=lambda k: intact[k][0])
    print()
    print("   Among rules that leave the outcome intact, removing %s"
          % worst)
    print("   closes the gap most: %.0f -> %.0f points."
          % (gap, intact[worst][0]))
    assert worst in ("R3", "R6")

    # --- step 3: read the chain for an affected ticket -------------
    print()
    print("=== step 3: read the explanation ===")
    for t in tickets:
        esc, fired, facts = decide(t)
        if t["team"] == "Platform" and esc == "no" and "R7" in fired:
            print("   ticket %s (%s), not escalated:"
                  % (t["id"], t["team"]))
            tree = explain(("escalate", "no"), facts)
            for line in render(tree):
                print("      " + line)
            break
    print()
    print("   R7 is a good rule. 'workaround' is recorded for 80% of")
    print("   Platform tickets and 15% of Payments tickets, because")
    print("   Platform has runbooks. A difference in DOCUMENTATION is")
    print("   being read as a difference in OPERATIONAL NEED.")

    # --- the competence boundary -----------------------------------
    print()
    print("=== what the rule base does NOT cover ===")
    import itertools
    attrs = {"severity": ["critical", "major", "minor"],
             "tier": ["platinum", "standard"],
             "workaround": ["yes", "no"],
             "owner_available": ["yes", "no"]}
    uncovered = []
    for combo in itertools.product(*attrs.values()):
        facts = dict(zip(attrs.keys(), combo))
        wm, _ = forward_chain(facts)
        if "escalate" not in wm:
            uncovered.append(facts)
    print("   %d of %d possible situations reach no escalate decision"
          % (len(uncovered), 3 * 2 * 2 * 2))
    for u in uncovered[:3]:
        print("      ", u)
    print("   A rule base can be asked this. Enumerate the inputs,")
    print("   find the gaps, and SAY SO -- that is a boundary the")
    print("   system can describe, and it takes an afternoon.")

    # --- the computational limit -----------------------------------
    print()
    print("=== and the limit no engineering removes ===")
    print("   b=4:  depth   nodes          at 1e6 nodes/sec")
    print("   " + "-" * 46)
    for d in (10, 15, 20, 25):
        n = 4 ** d
        secs = n / 1e6
        if secs < 90:
            t = "%.1f seconds" % secs
        elif secs < 5400:
            t = "%.1f minutes" % (secs / 60)
        elif secs < 86400 * 2:
            t = "%.1f hours" % (secs / 3600)
        else:
            t = "%.0f days" % (secs / 86400)
        print("   %11d %14d   %s" % (d, n, t))
    print()
    print("   A 100x faster machine buys three or four levels of")
    print("   depth and then stops. Structure is the only real")
    print("   answer, which is what Parts II to V were about.")
\end{lstlisting}
Then change \texttt{p\_workaround} so both teams document at the same rate, and
re-run. The gap collapses, and no rule changed --- which is the cleanest possible
demonstration that the disparity was never in the rule base at all.
\end{lab}

%---------------------------------------------------------------------
\begin{checkpoint}
\begin{tightnum}
  \item The escalation adviser's rules never mention the reporting team, and the
        escalation rates differ by $22$ percentage points. Explain how, and say why
        deleting an attribute would not have prevented it.
  \item The R3--R7 pathway was identified as carrying most of the gap, and both rules are sensible.
        State what the audit established and what it explicitly did not, and say
        whose decision the remainder is.
  \item A reviewer overrides the system's recommendation in $0.5\%$ of cases and
        reports this as evidence the system is working well. Give the alternative
        explanation and the measurement that distinguishes them.
\end{tightnum}
\end{checkpoint}

%---------------------------------------------------------------------
\begin{chaptersummary}
\begin{tight}
  \item Bias enters through \term{proxy variables}, uneven inputs, rule
        \term{interactions} and the choice of \term{measure}. None requires intent,
        and none is visible by reading the rules.
  \item \emph{Fairness through unawareness fails.} The escalation adviser never saw
        the team name and produced a $22$-point gap, because
        \texttt{workaround} is a proxy for having runbooks.
  \item An \term{audit} is three steps: measure outcomes by group, attribute the gap
        by \term{ablation}, and read the explanation chain for an affected case.
        Keep the sensitive attribute \emph{for the audit} and never for the
        decision.
  \item Ablation has a trap: the rule whose removal closes the gap completely was
        the one that produced the outcome at all, and equality by abolition is not
        fairness. \emph{Report the overall rate beside the gap} and discard any
        ablation that collapses it --- a specific instance of the general hazard of
        optimising a fairness metric instead of understanding a system.
  \item An \term{explanation} is an accountability artefact: it makes a decision
        \term{contestable} on a specific step and correctable by editing one rule.
        Those are properties the affected person needs.
  \item \term{Automation bias} makes reviewers stop reviewing, most strongly where
        the system is most often right. A near-zero override rate is a warning, not
        a success.
  \item \term{Brittleness} is being wrong with undiminished confidence just outside
        the region of competence, with no signal. A rule base cannot degrade
        gracefully --- but it \emph{can} enumerate the situations it does not cover
        and say so, which is a real and underused capability.
  \item Every method in Parts~II to~V is exponential, NP-hard or undecidable.
        Structure --- heuristics, propagation, pruning, conditional independence ---
        is the only answer; hardware buys three or four levels of depth and stops.
\end{tight}
\end{chaptersummary}

\begin{reviewq}
  \item \textbf{(MCQ)} Fairness through unawareness fails because of:
        (a)~measurement error (b)~proxy variables (c)~automation bias
        (d)~intractability.
  \item \textbf{(MCQ)} Ablation, in an audit, means: (a)~deleting sensitive
        attributes (b)~removing one rule at a time and re-measuring
        (c)~retraining (d)~sampling for review.
  \item \textbf{(MCQ)} A near-zero override rate on a decision-support system is
        best read as: (a)~the system is accurate (b)~a possible sign of automation
        bias (c)~evidence of good explanations (d)~a sign of brittleness.
  \item \textbf{(MCQ)} Brittleness means: (a)~being wrong sometimes (b)~being wrong
        with full confidence just outside the region of competence
        (c)~failing to terminate (d)~degrading smoothly.
  \item \textbf{(Short)} Explain why an explanation is an accountability property
        rather than a convenience, using the contestability of a specific step.
  \item \textbf{(Short)} The audit found the gap and did not say what to do. Argue
        that this is the right division of labour.
  \item \textbf{(Short)} Give the one thing a rule base can do about its competence
        boundary that a learned model cannot, and say what it costs to do.
  \item \textbf{(Applied)} Design an audit for a system that ranks job applications.
        State the groups you would measure, the outcome you would count, how you
        would attribute a gap to a component, and what you would present to the
        person who has to decide what to do about it.
\end{reviewq}
